% 单位阶跃函数（综述）
% license CCBYSA3
% type Wiki

本文根据 CC-BY-SA 协议转载翻译自维基百科\href{https://en.wikipedia.org/wiki/Heaviside_step_function}{相关文章}

海维赛德阶跃函数，又称单位阶跃函数，通常记作$H$或$\theta$有时也记作 $u$; 1或$\mathbb{1}$是一种以英国工程师奥利弗·海维赛德命名的阶跃函数。它在自变量为负时取值为0在自变量为正时取值为1。关于函数在零点处的取值$H(0)$存在不同的约定。该函数是阶跃函数类的一种典型代表，所有的阶跃函数都可以表示为它的平移的线性组合。

这一函数最初出现在运算微积分中，用于求解微分方程。在该背景下，它表示一个在某一时刻“打开”并无限期保持“打开”的信号。海维赛德在分析电报通信时发展了运算微积分，并将该函数表示为1。  
\subsection{表述}
取约定 $H(0) = 1$，海维赛德函数可以定义为：

\begin{itemize}
\item 分段函数：
$$
H(x) := 
\begin{cases} 
1, & x \geq 0 \\[6pt] 
0, & x < 0 
\end{cases}~
$$
\item 使用艾弗森括号记号：
$$
H(x) := [x \geq 0]~
$$
\item 作为示性函数：
$$
H(x) := \mathbf{1}_{x \geq 0} = \mathbf{1}_{\mathbb{R}_{+}}(x)~
$$
若取另一种约定 $H(0) = \tfrac{1}{2}$，则可以表示为：
\item 分段函数：
$$
H(x) := 
\begin{cases} 
1, & x > 0 \\[6pt] 
\tfrac{1}{2}, & x = 0 \\[6pt] 
0, & x < 0 
\end{cases}~
$$
\item 符号函数的线性变换：
$$
H(x) := \tfrac{1}{2}\left(\operatorname{sgn}x + 1\right)~
$$
\item 两个艾弗森括号的算术平均：
$$
H(x) := \frac{[x \geq 0] + [x > 0]}{2}~
$$
\item 双参数反正切函数（atan2）的单侧极限：
$$
H(x) := \lim_{\epsilon \to 0^{+}} \frac{\operatorname{atan2}(\epsilon, -x)}{\pi}~
$$
\item 作为超函数：
$$
H(x) := \left( 1 - \frac{1}{2\pi i}\log z, \; -\frac{1}{2\pi i}\log z \right)~
$$
或等价地：
$$
H(x) := \left( -\frac{\log-z}{2\pi i}, \; -\frac{\log-z}{2\pi i} \right),~
$$
其中 $\log z$ 表示复对数的主值。\\
其它在 $H(0)$ 处未定义的形式包括：
\item 分段函数：
$$
H(x) := 
\begin{cases} 
1, & x > 0 \\[6pt] 
0, & x < 0 
\end{cases}~
$$
\item 斜坡函数（ramp function）的导数：
$$
H(x) := \frac{d}{dx}\max\{x,0\}, \quad x \neq 0~
$$
\item 用绝对值函数表示：
$$
H(x) = \frac{x + |x|}{2x}~
$$
\end{itemize}
\subsection{与狄拉克 δ 函数的关系}
狄拉克 δ 函数是海维赛德函数的弱导数：
$$
\delta(x) = \frac{d}{dx}\, H(x),~
$$
因此，海维赛德函数可以被视为狄拉克 δ 函数的积分。有时写作：
$$
H(x) := \int_{-\infty}^{x} \delta(s)\, ds,~
$$
不过，对于 $x = 0$，这一展开式可能不成立（甚至无意义），这取决于采用哪种形式主义来定义涉及 $\delta$ 的积分。在这种语境下，海维赛德函数就是一个几乎必然为零的随机变量的累积分布函数（CDF）。（参见常量随机变量。）
\subsection{解析近似}
海维赛德阶跃函数的近似形式在生物化学和神经科学中很有用。例如，阶跃函数的逻辑近似，如Hill 方程和米氏–曼腾方程，可用于模拟细胞在化学信号刺激下的二元开关反应。

对于阶跃函数的平滑近似，可以使用逻辑函数：
$$
H(x) \approx \tfrac{1}{2} + \tfrac{1}{2}\tanh kx \;=\; \frac{1}{1 + e^{-2kx}},~
$$
其中，$k$ 越大，则在 $x = 0$ 处的过渡越陡峭。

如果取约定 $H(0) = \tfrac{1}{2}$，则在极限意义下成立：
$$
H(x) = \lim_{k \to \infty} \tfrac{1}{2}\bigl(1 + \tanh kx\bigr) 
= \lim_{k \to \infty} \frac{1}{1 + e^{-2kx}}.~
$$
\begin{figure}[ht]
\centering
\includegraphics[width=10cm]{./figures/0f53d68d0d0eaf06.png}
\caption{$\tfrac{1}{2} + \tfrac{1}{2}\tanh(kx) \;=\; \frac{1}{1+e^{-2kx}}$当 $k \to \infty$ 时，该表达式趋近于阶跃函数。} \label{fig_DWJYhs_1}
\end{figure}
阶跃函数还有许多其他平滑的解析近似形式。\(^\text{[1]}\)其中一些可能性是：
$$
\begin{aligned}
H(x) &= \lim_{k \to \infty} \left( \tfrac{1}{2} + \tfrac{1}{\pi} \arctan(kx) \right) \\[6pt]
H(x) &= \lim_{k \to \infty} \left( \tfrac{1}{2} + \tfrac{1}{2}\operatorname{erf}(kx) \right)
\end{aligned}~
$$
这些极限在逐点意义以及分布意义下都成立。然而，一般来说，逐点收敛并不必然意味着分布收敛，反之亦然。（不过，如果一个逐点收敛的函数序列中的所有函数都被某个“良好”的函数一致有界，那么其收敛也在分布意义下成立。）

更一般地说，任何连续概率分布的累积分布函数（CDF），只要它在零点附近集中，并且带有一个可控方差的参数，在方差趋于零的极限下，都可以作为阶跃函数的近似。例如，上述三种近似函数分别是以下常见概率分布的累积分布函数：逻辑分布、柯西分布和正态分布。
\subsection{非解析近似}
海维赛德阶跃函数的近似也可以通过平滑过渡函数来实现，例如当$1 \leq m \to \infty$时，定义函数：
$$
f(x) =
\begin{cases}
\tfrac{1}{2}\Bigl(1 + \tanh\!\Bigl(m \tfrac{2x}{1 - x^{2}}\Bigr)\Bigr), & |x| < 1 \\[8pt]
1, & x \geq 1 \\[6pt]
0, & x \leq -1
\end{cases}~
$$
\subsection{积分表示}
海维赛德阶跃函数常常可以用积分表示：
$$
\begin{aligned}
H(x) &= \lim_{\varepsilon \to 0^{+}} -\frac{1}{2\pi i} \int_{-\infty}^{\infty} \frac{1}{\tau + i\varepsilon}\, e^{-ix\tau}\, d\tau \\[8pt]
&= \lim_{\varepsilon \to 0^{+}} \;\frac{1}{2\pi i} \int_{-\infty}^{\infty} \frac{1}{\tau - i\varepsilon}\, e^{ix\tau}\, d\tau ,
\end{aligned}~
$$
其中，第二个表示式可以很容易地从第一个推出，因为阶跃函数是实函数，因此它等于其自身的复共轭。
\subsection{零点处的取值}
由于 $H$ 通常用于积分，而函数在单点的取值并不会影响其积分，因此在大多数情况下，$H(0)$ 取何值并不重要。实际上，当$H$被视为一个分布或 $L^{\infty}$ 空间（见 $L^p$ 空间）的元素时，讨论它在零点的具体取值甚至没有意义，因为这类对象只在“几乎处处”上定义。如果使用某种解析近似（如前面举的例子），那么通常会采用该近似在零点的极限值作为 $H(0)$。

不同的取值选择有不同的理由：
\begin{itemize}
\item 取 $H(0) = \tfrac{1}{2}$：这样图像具有旋转对称性；换句话说，$H - \tfrac{1}{2}$ 是一个奇函数。在这种情况下，与符号函数的关系对所有 $x$ 成立：
  $$
  H(x) = \tfrac{1}{2}\bigl(1 + \operatorname{sgn}(x)\bigr),~
  $$
  并且$\forall x,\; H(x) + H(-x) = 1$.
\item 取 $H(0) = 1$：这种取法保证 $H$ 是右连续的。例如，累积分布函数（CDF）通常被定义为右连续的，勒贝格–斯蒂尔切斯积分中的被积函数也是右连续的。在这种情况下，$H$ 就是一个闭半无限区间的示性函数：
  $$
  H(x) = \mathbf{1}_{[0,\infty)}(x).~
  $$
  对应的概率分布是退化分布。
\item 取 $H(0) = 0$：这种取法保证 $H$ 是左连续的。在这种情况下，$H$ 是开半无限区间的示性函数：
  $$
  H(x) = \mathbf{1}_{(0,\infty)}(x).~
  $$
\item 在泛函分析、优化与博弈论的语境下：通常将海维赛德函数定义为一个多值函数，以保持极限函数的连续性并确保某些解的存在。在这种情况下，海维赛德函数在零点返回一个区间：$H(0) = [0,1]$.
\end{itemize}
\subsection{离散形式}
单位阶跃函数的另一种形式，可以定义为一个函数$H : \mathbb{Z} \to \mathbb{R}$（即自变量是离散变量 $n$），定义为：
$$
H[n] =
\begin{cases}
0, & n < 0, \\[6pt]
1, & n \geq 0,
\end{cases}~
$$
或者使用半最大值约定：
$$
H[n] =
\begin{cases}
0, & n < 0, \\[6pt]
\tfrac{1}{2}, & n = 0, \\[6pt]
1, & n > 0,
\end{cases}~
$$
其中 $n$ 为整数。若 $n$ 是整数，则 $n < 0$ 意味着 $n \leq -1$，而 $n > 0$ 则意味着函数在 $n = 1$ 时取到 1。因此，在采用半最大值约定时，这个所谓的“阶跃函数”在区间 $[-1,1]$ 上表现得更像是“斜坡”（ramp-like）的行为，而不能严格意义上称为真正的阶跃函数。
